\paragraph{Splitting the selected spectra without changing detection.}
Let $j:A\hookrightarrow B$ be a finite injective local map of complete
discrete valuation rings, with uniformizers $t,u$. Write
$j(t)=vu^e$, where $v$ is a unit and $e\geq1$.
The finite torsion-free $A$-module $B$ is free. Choose a basis $b_a$
and its coordinate maps $p_a:B\to A$. Then
\[
 y=\sum_a j(p_a(y))b_a,
 \qquad p_a(j(x)y)=xp_a(y),
 \qquad p_a(u^{eN}B)\subset t^NA.
\]
The final inclusion proves continuity of every projection, including
in the ramified case.

Transport the representation, ordinary characters, local bases and
block partition along $j$. The original Selmer conditions are linear
in the cocycle and its local boundary witnesses: they prescribe block
support, continuity, the cocycle identity, transpose polarization,
inertial vanishing, and specified zero entries in
$w^{-1}(z-dY)w$. Thus pushing along $j$ or applying $p_a$ to the
cocycle and its witnesses preserves every condition. The same holds
for skew odd boundary primitives.

Consequently the given odd family remains independent modulo
boundaries after extension: project a relation to each $A$-coordinate.
It also continues to generate up to a power of the new uniformizer.
Indeed, if the projected identities use exponents $s_a$, take
$N=\max_a s_a$ and combine them with coefficients
\[
                 c_a=v^{-N}b_a j(t^{N-s_a}),
 \qquad
             \sum_a c_a j(t^{s_a}p_a(z))=u^{eN}z.
\]
This preserves all block and polarization conditions.

The same calculation preserves paired detection on the \emph{same}
finite set. Its coefficient module is
$W_B=W_A\otimes_A B$. Each projection commutes with the transported
twisted action, transpose and detecting trace:
\[
 p_a\!\left([\eta]\operatorname{tr}(z(h)\rho_{1,B}(h))\right)
    =[\eta]\operatorname{tr}(p_a(z(h))\rho_{1,A}(h)).
\]
Project a detected cocycle, use its polarized primitive over $A$, and
recombine with the displayed coefficients. This gives a polarized
primitive for $u^{eN}z$ in the actual module $W_B$.

For the selected separable monic characteristic polynomials, take a
finite common splitting field of $k((t))$. Its integral closure of
$k[[t]]$ is a finite complete discrete valuation ring, and all roots
belong to it. It has the form $k'[[u]]$: the roots of
$X^{|k'|}-X$ lift the residue field uniquely by Hensel's lemma, and
expansion in $u$ gives the isomorphism
\cite{stacks-selmer-dvr-extension,stacks-selmer-coefficient-field}.
For finitely many points, enlarge their new finite residue fields to
a common degree over $\mathbf F_\ell$. These constant-field
extensions are finite free, so detection and the odd families are
preserved throughout \cite{alpo-coefficient-extension}.

The transported point satisfies the same form, inertia and ordinary
identities. Its even-span condition follows by projecting an even
matrix and recombining the old span identities; inequalities between
ordinary characters persist by injectivity. The selected spectra
are split with distinct roots over coefficient rings having residue
fields of the same cardinality. Root differences need only be
nonzero; their valuations remain in the subsequent local estimate.

\begin{corollary}[Simultaneous auxiliary choice]
\label{cor:selmer-simultaneous-choice}
Given finitely many normalized points in odd characteristic, assume
the preceding open-subgroup finiteness condition, the nontrivial
residual cyclotomic character, and spectral density in the original
cyclotomic kernel at each point. There are finite injective local
extensions of their coefficient rings, with residue fields of one
common cardinality, and one finite set of elements of that kernel
such that all selected characteristic polynomials split with distinct
roots. Each extended point has an independent odd Selmer family
spanning up to powers of its uniformizer, and the same finite set
detects its paired cocycles. The original local bases, block
partitions, ordinary characters and their prescribed relations are
preserved.
\end{corollary}

\begin{proof}
Construct the odd families from the finite Selmer modules as above.
Apply integral paired detection and finite simultaneous selection,
then the splitting and coefficient-extension argument. Every step
uses the original cyclotomic kernel and the same finite set.
\end{proof}
